\chapter{核能：裂变、聚变、$E=mc^2$}\label{ch:g12:nuclear-energy}

一粒铅笔橡皮大小的燃料芯块，可供一户家庭的
用电一年；同样的电量若改用煤来供应，则需三吨煤。头顶上，太阳
以自身消亡换取日光——每秒四百万吨。
本章将揭开这两笔交易背后的账本：质量本身就是
\omterm{def:g10:energy-conservation:energy}{能量}，而\omterm{def:g11:fundamental-interactions:ladder}{原子核}（见\cref{ch:g11:nucleus-radioactivity}）正是兑现它的地方。

\section{质量即能量}

\begin{theorem}[质能等价]\label{thm:g12:nuclear-energy:emc2}\index{质能等价}
任何质量为 $m$ 的物体，仅因其存在，便具有\omterm{def:g10:energy-conservation:energy}{能量}
\[
E = m c^2, \qquad c = \qty{3.00e8}{m/s}.
\]
每当一个系统损失质量 $\Delta m$ 时，便释放出\omterm{def:g10:energy-conservation:energy}{能量}
$\Delta m\, c^2$——无论机制如何。
\end{theorem}
\admitted

\begin{remark}[爱因斯坦，1905]\label{rem:g12:nuclear-energy:einstein}
这是爱因斯坦关系式，在本阶段不加证明地承认；其来源
属于本学年末的狭义相对论一章
（见\cref{ch:g12:special-relativity}）。注意其变化率：$c^2 = \qty{9.00e16}{J/kg}$
——质量上的微小舍入误差，便是一笔可观的\omterm{def:g10:energy-conservation:energy}{能量}。
\end{remark}

\begin{definition}[电子伏特与原子质量单位]\label{def:g12:nuclear-energy:units}
两个专为\omterm{def:g11:fundamental-interactions:ladder}{原子核}设计的\omterm{def:g10:orders-of-magnitude:unit}{单位}。
其一是\emph{电子伏特}\index{电子伏特}，即电子
穿过 \qty{1}{V} 电势差（见\cref{ch:g11:circuits-and-power}）时所获得的\omterm{def:g10:energy-conservation:energy}{能量}：
$\qty{1}{eV} = \qty{1.60e-19}{J}$, $\qty{1}{MeV} = \qty{1.60e-13}{J}$ ---
化学键的\omterm{def:g10:energy-conservation:energy}{能量}只有几个 \unit{eV}，核反应则是兆电子\omterm{def:g11:circuits-and-power:voltage}{伏特}。
另一个是\emph{原子质量单位}\index{原子质量单位}，
即碳-12 \omterm{def:g11:fundamental-interactions:ladder}{原子}质量的十二分之一：$\qty{1}{u} = \qty{1.66054e-27}{kg}$，
其\omterm{def:g10:energy-conservation:energy}{能量}当量为 \qty{931.5}{MeV}——$\qty{1}{u} = \qty{931.5}{MeV}/c^2$。
\end{definition}

\begin{example}[沉睡的财富]\label{ex:g12:nuclear-energy:fortune}
任何一\omterm{def:g10:orders-of-magnitude:unit}{千克}物质，若完全转化，即为 $\qty{9.00e16}{J}$——相当于一座
大型（\qty{1}{GW}）发电站三年的产出。人们直到
1905 年才注意到，因为寻常物理几乎不涉及质量：燃烧一\omterm{def:g10:orders-of-magnitude:unit}{千克}煤
（\qty{3.0e7}{J}）使其质量减轻三分之一微克。后文将看到，\omterm{def:g11:fundamental-interactions:ladder}{原子核}
的变化幅度为千分之一：仍然很小，却是化学过程的一百万倍。
\end{example}

\section{质量亏损}

\begin{definition}[质量亏损与结合能]\label{def:g12:nuclear-energy:defect}
称量一个\omterm{def:g11:fundamental-interactions:ladder}{原子核} ${}^{A}_{Z}\mathrm{X}$（质量为 $m$），再分别称量其 $Z$ 个质子
与 $A - Z$ 个中子：各部分的总和比整体\emph{更重}。
其差值为
\[
\Delta m = Z\,m_p + (A - Z)\,m_n - m > 0,
\qquad m_p = \qty{1.00728}{u},\; m_n = \qty{1.00866}{u},
\]
即\emph{质量亏损}\index{质量亏损}；
\emph{结合能}\index{结合能} $E_b = \Delta m\,c^2$ 是
将\omterm{def:g11:fundamental-interactions:ladder}{原子核}拆散为自由\omterm{def:g11:fundamental-interactions:ladder}{核子}所需的\omterm{def:g10:energy-conservation:energy}{能量}——同理，
组装时所释放的\omterm{def:g10:energy-conservation:energy}{能量}。束缚态的\omterm{def:g11:fundamental-interactions:ladder}{原子核}位于其组分\emph{之下}。
\end{definition}

\begin{omfigure}
\begin{tikzpicture}[font=\small]
  \draw[thick] (0,2.1) -- (2.7,2.1) node[right] {$Z$ 个质子 $+$ $(A-Z)$ 个中子，自由};
  \foreach \x in {0.3,0.9,1.5,2.1} \fill[omThm!80] (\x,2.4) circle (2.4pt);
  \foreach \x in {0.6,1.2,1.8,2.4} \fill[omExo] (\x,2.4) circle (2.4pt);
  \draw[thick, omDef] (0,0) -- (2.7,0) node[right] {原子核，质量 $m$};
  \draw[-Stealth, thick] (1.35,2.1) -- (1.35,0) node[midway, right] {$E_b = \Delta m\,c^2$};
\end{tikzpicture}

{\small \omterm{def:g10:energy-conservation:energy}{能量}阶梯示意：束缚态\omterm{def:g11:fundamental-interactions:ladder}{原子核}比其分离的\omterm{def:g11:fundamental-interactions:ladder}{核子}低 $E_b$，
\omterm{def:g11:fundamental-interactions:ladder}{核子}组装成\omterm{def:g11:fundamental-interactions:ladder}{原子核}时释放 $E_b$，拆解则需要付出这份\omterm{def:g10:energy-conservation:energy}{能量}。}
\end{omfigure}

\begin{example}[氦-4]\label{ex:g12:nuclear-energy:helium}
${}^{4}_{2}\mathrm{He}$ \omterm{def:g11:fundamental-interactions:ladder}{原子核}的质量为 $m = \qty{4.00151}{u}$；其组分为：
$2 \times 1.00728 + 2 \times 1.00866 = \qty{4.03188}{u}$。
故 $\Delta m = \qty{0.03037}{u}$，则 $E_b = 0.03037 \times 931.5 \approx
\qty{28.3}{MeV}$：该\omterm{def:g11:fundamental-interactions:ladder}{原子核}比其组分轻 $0.75\%$——这与
\cref{ex:g12:nuclear-energy:fortune}中的“千分之一”相呼应。
\end{example}

\begin{method}[核反应的能量平衡]\label{met:g12:nuclear-energy:balance}
\begin{enumerate}
  \item 配平方程：$A$ 与 $Z$ 的总和守恒（见\cref{ch:g11:nucleus-radioactivity}）。
  \item 分别求反应前、后的质量之和，\omterm{def:g10:orders-of-magnitude:unit}{单位}为 \unit{u}——保留全部五位小数。
  \item $\Delta m = m_{\text{反应前}} - m_{\text{反应后}}$，然后
        $E = \Delta m\,(\text{以 u 计}) \times \qty{931.5}{MeV}$，以
        生成物\omterm{def:g11:mechanical-energy:kinetic}{动能}的形式带走；若为负值，则该反应需外界付出\omterm{def:g10:energy-conservation:energy}{能量}。
\end{enumerate}
\end{method}

\section{驱动宇宙的曲线}

\begin{definition}[比结合能]\label{def:g12:nuclear-energy:pernucleon}
结合程度的公平比较，依据的是
\emph{比结合能}\index{比结合能}，即 $E_b/A$，是拆出
一个\omterm{def:g11:fundamental-interactions:ladder}{核子}的平均代价：对氦-4 而言为 $28.3/4 \approx \qty{7.1}{MeV}$。
\end{definition}

\begin{proposition}[铁峰与两条下行之路]\label{prop:g12:nuclear-energy:ironpeak}\index{铁峰}
以 $A$ 为横坐标作图，\omterm{def:g12:nuclear-energy:pernucleon}{比结合能}在经过轻核时陡然上升，
在铁（${}^{56}_{26}\mathrm{Fe}$）附近达到峰值，约为
\qty{8.8}{MeV}（每个\omterm{def:g11:fundamental-interactions:ladder}{核子}），随后缓慢下降，在
铀处约为 \qty{7.6}{MeV}：铁是自然界中结合最紧密的\omterm{def:g11:fundamental-interactions:ladder}{原子核}。
一个反应释放\omterm{def:g10:energy-conservation:energy}{能量}，当且仅当其生成物在曲线上\emph{位置更高}；因此两种
截然相反的策略都有所得，且都走向铁——
使最重的\omterm{def:g11:fundamental-interactions:ladder}{原子核}\emph{\omterm{def:g12:nuclear-energy:fission}{裂变}}，使最轻的\omterm{def:g11:fundamental-interactions:ladder}{原子核}\emph{\omterm{def:g12:nuclear-energy:fusion}{聚变}}。
\end{proposition}
\admitted

\begin{remark}[为何有峰]\label{rem:g12:nuclear-energy:whypeak}
\cref{ch:g11:fundamental-interactions}中的拉锯：强“胶水”
只束缚相邻的\omterm{def:g11:fundamental-interactions:ladder}{核子}，\omterm{def:g11:fundamental-interactions:charge}{库仑}斥力作用于整个\omterm{def:g11:fundamental-interactions:ladder}{原子核}——小核
多为表面，大核多受斥力，铁则是二者的折中。
严谨的核算见本科分册。
\end{remark}

\begin{omfigure}
\begin{tikzpicture}
\begin{axis}[omaxis, width=9.2cm, height=5.6cm,
    xlabel={质量数 $A$}, ylabel={$E_b/A$ (\unit{MeV})},
    xmin=0, xmax=248, ymin=0, ymax=10.4, ytick={0,2,4,6,8}, xtick={4,56,120,180,235}]
  \addplot[omDef, thick, smooth] coordinates {(1,0) (2,1.1) (3,2.8) (6,5.3) (9,6.5)
    (12,7.7) (20,8.0) (40,8.6) (56,8.8) (90,8.7) (120,8.5) (150,8.3) (190,7.9) (235,7.6)};
  \addplot[only marks, mark=*, mark size=1.6pt, omThm] coordinates {(4,7.1) (56,8.8) (235,7.6)};
  \node[font=\small, right] at (axis cs:8,7.1) {$^{4}$He};
  \node[font=\small, above] at (axis cs:56,8.9) {$^{56}$Fe};
  \node[font=\small, above] at (axis cs:235,7.75) {$^{235}$U};
  \draw[-Stealth, omProp, thick] (axis cs:16,3.4) -- (axis cs:44,7.6)
    node[midway, below right, font=\small] {聚变};
  \draw[-Stealth, omProp, thick] (axis cs:222,6.9) -- (axis cs:130,7.8)
    node[midway, below, font=\small] {裂变};
\end{axis}
\end{tikzpicture}

{\small 每个\omterm{def:g11:fundamental-interactions:ladder}{核子}的\omterm{def:g12:nuclear-energy:defect}{结合能}：陡升、铁峰、缓降——
\omterm{def:g12:nuclear-energy:fusion}{聚变}自左侧攀升，\omterm{def:g12:nuclear-energy:fission}{裂变}自右侧下行，两者都释放\omterm{def:g10:energy-conservation:energy}{能量}。}
\end{omfigure}

\section{裂变与链式反应}

\begin{definition}[裂变]\label{def:g12:nuclear-energy:fission}
\emph{裂变}\index{裂变}是重\omterm{def:g11:fundamental-interactions:ladder}{原子核}分裂为两个
中等质量的碎片。铀-235 是\emph{易裂变的}\index{易裂变的}：一个慢
中子被吸收后使\omterm{def:g11:fundamental-interactions:ladder}{原子核}激发得十分剧烈，以至于一分为二，
并放出两个或三个新的中子——
${}^{1}_{0}\mathrm{n} + {}^{235}_{92}\mathrm{U} \to \text{两个碎片} +
2\text{--}3\;{}^{1}_{0}\mathrm{n}$，放出的中子数随每次裂变而异。
\end{definition}

\begin{example}[称量一次裂变]\label{ex:g12:nuclear-energy:fissionbalance}
依\cref{met:g12:nuclear-energy:balance}，一条常见途径为：
\[
{}^{1}_{0}\mathrm{n} + {}^{235}_{92}\mathrm{U} \longrightarrow
{}^{94}_{38}\mathrm{Sr} + {}^{140}_{54}\mathrm{Xe} + 2\,{}^{1}_{0}\mathrm{n}.
\]
\omterm{def:g12:nuclear-energy:fission}{裂变}前：$1.00866 + 234.99350 = \qty{236.00216}{u}$。\omterm{def:g12:nuclear-energy:fission}{裂变}后：
$93.89450 + 139.89200 + 2 \times 1.00866 = \qty{235.80382}{u}$。由此得
$\Delta m = \qty{0.19834}{u}$，$E \approx \qty{185}{MeV}$；计入
\omterm{def:g11:nucleus-radioactivity:radioactivity}{放射性}碎片的后续衰变（见\cref{ch:g12:radioactive-decay}），
每次\omterm{def:g12:nuclear-energy:fission}{裂变}约合 \qty{200}{MeV}——几个 \unit{eV} 的\omterm{def:g10:energy-conservation:energy}{能量}只够
燃烧一个煤\omterm{def:g11:fundamental-interactions:ladder}{原子}。
\end{example}

\begin{definition}[链式反应与临界质量]\label{def:g12:nuclear-energy:chain}
每次\omterm{def:g12:nuclear-energy:fission}{裂变}由一个中子引发，并释放两到三个中子：块体足够大时，
这些中子撞击其他\omterm{def:g11:fundamental-interactions:ladder}{原子核}，反应便自我维持——即
\emph{链式反应}\index{链式反应}。在小块体中，多数中子
先从表面逸出；维持链式反应的最小质量
即\emph{临界质量}\index{临界质量}——低于它，链式反应
便会熄灭；高于它，每一代都成倍增长：这正是炸弹的原理。
\end{definition}

\begin{omfigure}
\begin{tikzpicture}[font=\small, scale=0.9]
  \draw[-Stealth, thick] (-1.4,0) -- (-0.4,0) node[pos=0.15, above] {n};
  \fill[omDef!55] (0,0) circle (11pt); \node at (0,0) {$^{235}$U};
  \fill[omExo] (0.75,0.85) circle (4.5pt) node[above=3pt] {碎片};
  \fill[omExo] (0.75,-0.85) circle (4.5pt);
  \foreach \y in {1.8,0,-1.8} {\draw[-Stealth] (0.5,0) -- (2.5,\y);
    \fill[omDef!55] (3,\y) circle (11pt); \node at (3,\y) {$^{235}$U};
    \foreach \a in {28,0,-28} \draw[-Stealth] (3.5,\y) -- ++(\a:1.1);}
\end{tikzpicture}

{\small \omterm{def:g12:nuclear-energy:chain}{链式反应}：一个中子裂开一个\omterm{def:g11:fundamental-interactions:ladder}{原子核}，释放出的三个中子
又去裂开另外三个\omterm{def:g11:fundamental-interactions:ladder}{原子核}——反应堆只允许其中恰好一个继续。}
\end{omfigure}

\begin{definition}[驯服的反应堆]\label{def:g12:nuclear-energy:reactor}
动力反应堆将\omterm{def:g12:nuclear-energy:chain}{链式反应}控制为每次\omterm{def:g12:nuclear-energy:fission}{裂变}恰好保留一个中子。燃料：
经适度富集的铀-235 燃料芯块。
\emph{慢化剂}\index{慢化剂}（通常为普通水）使快速的\omterm{def:g12:nuclear-energy:fission}{裂变}中子
通过碰撞减速：慢中子远比快中子更易引发
下一次\omterm{def:g12:nuclear-energy:fission}{裂变}；\emph{控制棒}\index{控制棒}（硼、镉）
吞噬中子——插入则\omterm{def:g12:nuclear-energy:chain}{链式反应}熄灭，抽出则反应加快。
碎片的剧烈运动转化为热，再转化为蒸汽，最后转化为\omterm{def:g10:energy-conservation:forms}{电能}
（见\cref{ch:g10:energy-conservation}）。
\end{definition}

\begin{example}[胡椒粒对卡车]\label{ex:g12:nuclear-energy:gram}
一克铀-235 含有
$N = \num{1.0e-3} / (235 \times \num{1.66054e-27}) \approx \num{2.56e21}$
个\omterm{def:g11:fundamental-interactions:ladder}{原子核}；每个以 $\qty{200}{MeV} = \qty{3.2e-11}{J}$ 计，全部\omterm{def:g12:nuclear-energy:fission}{裂变}
约释放 \qty{8.2e10}{J}——相当于 $\num{8.2e10} /
\num{3.0e7} \approx \qty{2.7e3}{kg}$ 煤的\omterm{def:g10:energy-conservation:forms}{化学能}：一粒胡椒对三吨。
\end{example}

\section{聚变：太阳自身的火焰}

\begin{definition}[聚变]\label{def:g12:nuclear-energy:fusion}
\emph{聚变}\index{聚变}是两个轻\omterm{def:g11:fundamental-interactions:ladder}{原子核}合并为一个更重、
结合更紧密的\omterm{def:g11:fundamental-interactions:ladder}{原子核}。最易实现的反应使氘与
氚结合，它们是氢的重\omterm{def:g11:nucleus-radioactivity:notation}{同位素}（见\cref{ch:g11:nucleus-radioactivity}）：
\[
{}^{2}_{1}\mathrm{H} + {}^{3}_{1}\mathrm{H} \longrightarrow
{}^{4}_{2}\mathrm{He} + {}^{1}_{0}\mathrm{n}.
\]
\end{definition}

\begin{example}[D--T，称量]\label{ex:g12:nuclear-energy:dt}
反应前：$2.01355 + 3.01550 = \qty{5.02905}{u}$；反应后：
$4.00151 + 1.00866 = \qty{5.01017}{u}$；于是 $\Delta m = \qty{0.01888}{u}$，
$E \approx \qty{17.6}{MeV}$，其中五分之四落在中子上——少于一次
\omterm{def:g12:nuclear-energy:fission}{裂变}，但其\omterm{def:g10:energy-conservation:energy}{能量}来自五个\omterm{def:g11:fundamental-interactions:ladder}{核子}而非 $236$ 个，约合 \qty{3.4e14}{J}
每\omterm{def:g10:orders-of-magnitude:unit}{千克}燃料，是\omterm{def:g12:nuclear-energy:fission}{裂变}的四倍。氦-4
（见\cref{ex:g12:nuclear-energy:helium}）仍是“灰烬”。
\end{example}

\begin{definition}[库仑势垒]\label{def:g12:nuclear-energy:barrier}
两个\omterm{def:g11:fundamental-interactions:ladder}{原子核}都带正电：要接触，它们必须先爬上
它们之间电斥力筑成的山坡（\cref{ch:g11:electric-gravitational-fields}）——
\emph{库仑势垒}\index{库仑势垒}。只有接近 \qty{1e8}{K} 时，
碰撞才能提供越过它的\omterm{def:g10:energy-conservation:energy}{能量}；此时燃料是
裸核与电子构成的\emph{等离子体}\index{等离子体}。\omterm{def:g12:nuclear-energy:fission}{裂变}则无此代价：
其引发者中子不带电，在低温下也能进行。
\end{definition}

\begin{example}[太阳靠它发光]\label{ex:g12:nuclear-energy:sun}
太阳核心温度为 $\qty{1.5e7}{K}$，通过逐步\omterm{def:g12:nuclear-energy:fusion}{聚变}普通氢——
\emph{质子–质子链}，其净效应为
$4\,{}^{1}_{1}\mathrm{H} \to {}^{4}_{2}\mathrm{He} + 2\,{}^{0}_{+1}\mathrm{e}
+ \text{辐射}$，每生成一个氦约释放 \qty{26}{MeV}，$0.7\%$ 的
氢质量转化为\omterm{def:g10:energy-conservation:energy}{能量}。我们在
\cref{ch:g10:light-spectra}中读取\omterm{def:g10:light-spectra:white}{光谱}的太阳光，即是这部分质量，晚八分钟抵达；每颗恒星
都沿着曲线向铁的方向“行走”而发光。
\end{example}

\begin{remark}[ITER：把恒星装进瓶中]\label{rem:g12:nuclear-energy:iter}
任何固体都无法在 \qty{1.5e8}{K} 下容纳\omterm{def:g12:nuclear-energy:barrier}{等离子体}；\emph{托卡马克}用
\omterm{def:g11:magnetic-fields:bfield}{磁场}（\cref{ch:g11:magnetic-fields}）将其围住，形如甜甜圈的
\omterm{def:g11:electric-gravitational-fields:lines}{场线}区域，带电粒子沿其螺旋运动而不触及器壁；ITER 旨在释放
加热\omterm{def:g11:work-of-force:power}{功率}十倍的\omterm{def:g10:energy-conservation:energy}{能量}，是当前的一步。一旦破坏
约束，\omterm{def:g12:nuclear-energy:barrier}{等离子体}便会在数秒内降温并停止：没有\omterm{def:g12:nuclear-energy:chain}{链式反应}，没有
\omterm{def:g12:nuclear-energy:chain}{临界质量}——难点不在于阻止火焰，而在于让它保持燃烧。
\end{remark}

\begin{example}[每千克的能量阶梯]\label{ex:g12:nuclear-energy:ladder}
每\omterm{def:g10:orders-of-magnitude:unit}{千克}燃料的\omterm{def:g10:energy-conservation:energy}{能量}：煤，\qty{3.0e7}{J}；铀-235 \omterm{def:g12:nuclear-energy:fission}{裂变}，
\qty{8.2e13}{J}；氘–氚\omterm{def:g12:nuclear-energy:fusion}{聚变}，\qty{3.4e14}{J}；全部
质量转化（$E = mc^2$），\qty{9.0e16}{J}。化学与\omterm{def:g11:fundamental-interactions:ladder}{原子核}之间
相差六个\omterm{def:g10:orders-of-magnitude:oom}{数量级}——这正是反应堆以卡车换料、
燃煤电厂以整列火车运煤的原因。
\end{example}

\section{练习}

\begin{exercise}[$\star$]\label{exo:g12:nuclear-energy:1}
将下列各量换算为\omterm{def:g11:work-of-force:work}{焦耳}：(a) \qty{1}{eV} 与 \qty{1}{MeV}；(b) 一次\omterm{def:g12:nuclear-energy:fission}{裂变}的
\qty{200}{MeV}；(c) 由 $\qty{1}{u} = \qty{1.66054e-27}{kg}$ 验证其
\omterm{def:g10:energy-conservation:energy}{能量}当量接近 \qty{931.5}{MeV}。
\end{exercise}

\begin{exercise}[$\star$]\label{exo:g12:nuclear-energy:2}
一块方糖的质量为 \qty{6.0}{g}，求其完整的 $E = mc^2$ 值。
一户家庭每天耗电 $\qty{10}{kW.h} = \qty{3.6e7}{J}$。若将这块方糖
全部转化，可供其用电多少年？
\end{exercise}

\begin{exercise}[$\star$]\label{exo:g12:nuclear-energy:3}
碳-12 \omterm{def:g11:fundamental-interactions:ladder}{原子核}的质量为 \qty{11.99671}{u}，求其\omterm{def:g12:nuclear-energy:defect}{质量亏损}，
其\omterm{def:g12:nuclear-energy:defect}{结合能}（以 \unit{MeV} 计），以及其\omterm{def:g12:nuclear-energy:pernucleon}{比结合能}
（$m_p = \qty{1.00728}{u}$，$m_n = \qty{1.00866}{u}$）。
\end{exercise}

\begin{exercise}[$\star$]\label{exo:g12:nuclear-energy:4}
补全下列各\omterm{def:g12:nuclear-energy:fission}{裂变}，给出 $A$、$Z$ 及元素（$Z = 36$：氪，
$Z = 52$：碲）：
(a) ${}^{1}_{0}\mathrm{n} + {}^{235}_{92}\mathrm{U} \to
{}^{144}_{56}\mathrm{Ba} + {}^{A}_{Z}\mathrm{X} + 3\,{}^{1}_{0}\mathrm{n}$;
(b) ${}^{1}_{0}\mathrm{n} + {}^{235}_{92}\mathrm{U} \to
{}^{97}_{40}\mathrm{Zr} + {}^{A}_{Z}\mathrm{Y} + 2\,{}^{1}_{0}\mathrm{n}$.
\end{exercise}

\begin{exercise}[$\star$]\label{exo:g12:nuclear-energy:5}
读取\cref{prop:g12:nuclear-energy:ironpeak}中的曲线，估算下列各核的 $E_b/A$：
氦-4、碳-12、铁-56 与铀-235。哪一种结合最紧密？
哪一种可由\omterm{def:g12:nuclear-energy:fission}{裂变}释放\omterm{def:g10:energy-conservation:energy}{能量}，哪一种可由\omterm{def:g12:nuclear-energy:fusion}{聚变}释放\omterm{def:g10:energy-conservation:energy}{能量}？
\end{exercise}

\begin{exercise}[$\star\star$]\label{exo:g12:nuclear-energy:6}
氘核 ${}^{2}_{1}\mathrm{H}$ 的质量为 \qty{2.01355}{u}，求其
\omterm{def:g12:nuclear-energy:defect}{结合能}（以 \unit{MeV} 计）。一个化学键束缚的\omterm{def:g10:energy-conservation:energy}{能量}约为几个 \unit{eV}——
核“键”比它强多少倍？
\end{exercise}

\begin{exercise}[$\star\star$]\label{exo:g12:nuclear-energy:7}
核算下述\omterm{def:g12:nuclear-energy:fission}{裂变}前后的质量：${}^{1}_{0}\mathrm{n} + {}^{235}_{92}\mathrm{U} \to
{}^{144}_{56}\mathrm{Ba} + {}^{89}_{36}\mathrm{Kr} + 3\,{}^{1}_{0}\mathrm{n}$:
质量分别为 \qty{234.99350}{u}（U）、\qty{143.89220}{u}（Ba）、\qty{88.89790}{u}
（Kr）、\qty{1.00866}{u}（n）。释放的\omterm{def:g10:energy-conservation:energy}{能量}以 \unit{MeV} 和\omterm{def:g11:work-of-force:work}{焦耳}计，各为多少？
\end{exercise}

\begin{exercise}[$\star\star$]\label{exo:g12:nuclear-energy:8}
另一\omterm{def:g12:nuclear-energy:fusion}{聚变}反应：${}^{2}_{1}\mathrm{H} + {}^{2}_{1}\mathrm{H} \to
{}^{3}_{2}\mathrm{He} + {}^{1}_{0}\mathrm{n}$，其中
$m({}^{3}_{2}\mathrm{He}) = \qty{3.01493}{u}$。求释放的\omterm{def:g10:energy-conservation:energy}{能量}。
为何远少于氘–氚反应？（考虑各生成物在曲线上的位置。）
\end{exercise}

\begin{exercise}[$\star\star$]\label{exo:g12:nuclear-energy:9}
各用几行解释：(a) 为何\omterm{def:g12:nuclear-energy:fission}{裂变}无需加热，而\omterm{def:g12:nuclear-energy:fusion}{聚变}
需要约 \qty{1e8}{K}；(b) 太阳如何在“仅”\qty{1.5e7}{K} 下运作
（以人群的规模及其中最快的成员作比喻）；(c) 为何\omterm{def:g12:nuclear-energy:fusion}{聚变}电站
无法像\omterm{def:g12:nuclear-energy:chain}{链式反应}那样失控。
\end{exercise}

\begin{exercise}[$\star\star$]\label{exo:g12:nuclear-energy:10}
太阳的净反应将四个质子（$m_p = \qty{1.00728}{u}$）转化为一个
氦-4 \omterm{def:g11:fundamental-interactions:ladder}{原子核}（\qty{4.00151}{u}）：求每生成一个氦所释放的\omterm{def:g10:energy-conservation:energy}{能量}，以
\unit{MeV} 计，以及初始质量中被转化的比例。
\end{exercise}

\begin{exercise}[$\star\star$]\label{exo:g12:nuclear-energy:11}
在以水为\omterm{def:g12:nuclear-energy:reactor}{慢化剂}的反应堆中：(a) \omterm{def:g12:nuclear-energy:reactor}{慢化剂}起什么作用，为何
\omterm{def:g12:nuclear-energy:chain}{链式反应}需要它？(b) \omterm{def:g12:nuclear-energy:reactor}{控制棒}呢？(c) 冷却水\emph{即是}\omterm{def:g12:nuclear-energy:reactor}{慢化剂}：为何失去它往往
会抑制\omterm{def:g12:nuclear-energy:chain}{链式反应}，而非使其加速？
\end{exercise}

\begin{exercise}[$\star\star\star$]\label{exo:g12:nuclear-energy:12}
一条不受控制的\omterm{def:g12:nuclear-energy:chain}{链式反应}每一代都会翻倍。从一次\omterm{def:g12:nuclear-energy:fission}{裂变}出发，
大约需要多少代，一克铀-235 中的 $\num{2.56e21}$ 个\omterm{def:g11:fundamental-interactions:ladder}{原子核}才会全部
\omterm{def:g12:nuclear-energy:fission}{裂变}完毕（$2^{10} \approx 10^{3}$）？若每代约需 \qty{10}{ns}，
这一过程共需多久？由此可见，为何临界状态要受到如此严格的控制？
\end{exercise}

\begin{exercise}[$\star\star\star$]\label{exo:g12:nuclear-energy:13}
我们能否通过使铁\omterm{def:g12:nuclear-energy:fission}{裂变}来获取\omterm{def:g10:energy-conservation:energy}{能量}？考察
${}^{56}_{26}\mathrm{Fe} \to 2\,{}^{28}_{14}\mathrm{Si}$ 这一反应，其中
$m(\mathrm{Fe}) = \qty{55.92066}{u}$，$m(\mathrm{Si}) = \qty{27.96924}{u}$。
计算 $\Delta m$，并结合曲线得出结论，说明为何铁是
核燃烧的灰烬：燃料一去不复返。
\end{exercise}

\begin{exercise}[$\star\star\star$]\label{exo:g12:nuclear-energy:14}
由 $m({}^{56}_{26}\mathrm{Fe}) = \qty{55.92066}{u}$ 与
$m({}^{235}_{92}\mathrm{U}) = \qty{234.99350}{u}$，分别计算二者的 $E_b/A$。
（$m_p = \qty{1.00728}{u}$，$m_n = \qty{1.00866}{u}$）。\omterm{def:g12:nuclear-energy:fission}{裂变}碎片的
\omterm{def:g12:nuclear-energy:pernucleon}{比结合能}约为每个\omterm{def:g11:fundamental-interactions:ladder}{核子} \qty{8.5}{MeV}：请根据这两个能级，\emph{估算}一次\omterm{def:g12:nuclear-energy:fission}{裂变}所释放的\omterm{def:g10:energy-conservation:energy}{能量}，
并与\cref{ex:g12:nuclear-energy:fissionbalance}作比较。
\end{exercise}

\begin{exercise}[$\star\star\star$]\label{exo:g12:nuclear-energy:15}
ITER 的目标是提供 \qty{500}{MW} 的\omterm{def:g12:nuclear-energy:fusion}{聚变}\omterm{def:g11:work-of-force:power}{功率}。每次 D--T 反应放出 \qty{17.6}{MeV}
（消耗 \qty{5.029}{u} 的燃料），求每秒发生的反应次数，
再求每天消耗的燃料质量。一座 \qty{500}{MW} 的燃煤电厂约烧掉
每天 \qty{1.4e6}{kg}：请比较两者。
\end{exercise}

\section{问题：为一座城市供电}

\begin{problem}[{周末问题——为一座城市供电：一吉瓦的稳定功率，
供五十万人使用，可由四种方式获取——铀芯块、成列火车运来的煤、
一池海水、太阳自身的物质——四者背后都隐藏着同样一千克的
消失质量}]
\label{pb:g12:nuclear-energy:1}
一座拥有 $500\,000$ 名居民的城市，其持续电\omterm{def:g11:work-of-force:power}{功率}为
$P = \qty{1.0}{GW}$；无论使用何种燃料，其发电厂都将\omterm{def:g10:energy-conservation:forms}{热能}转化为
\omterm{def:g10:energy-conservation:forms}{电能}，\omterm{def:g10:energy-conservation:efficiency}{效率}为 $33\%$（见\cref{ch:g10:energy-conservation}）。
数据：一次铀-235 \omterm{def:g12:nuclear-energy:fission}{裂变}共放出 \qty{200}{MeV}；煤的热值为 \qty{3.0e7}{J/kg}；
一次 D--T \omterm{def:g12:nuclear-energy:fusion}{聚变}由 \qty{5.029}{u} 的燃料放出 \qty{17.6}{MeV}；
$\qty{1}{u} = \qty{1.66054e-27}{kg}$；$\qty{1}{eV} = \qty{1.60e-19}{J}$；一
年 $= \qty{3.16e7}{s}$；太阳：\qty{3.9e26}{W}，\qty{2.0e30}{kg}；一立方
米海水中约含 \qty{33}{g} 氘。

\medskip
\noindent\textbf{第一部分——\omterm{def:g12:nuclear-energy:fission}{裂变}电厂。}
\begin{enumerate}
  \item 反应堆必须提供多大的热\omterm{def:g11:work-of-force:power}{功率}？
  \item 用\omterm{def:g11:work-of-force:work}{焦耳}表示一次\omterm{def:g12:nuclear-energy:fission}{裂变}的\omterm{def:g10:energy-conservation:energy}{能量}。
  \item 每秒需要发生多少次\omterm{def:g12:nuclear-energy:fission}{裂变}，才能维持这座城市的用电？
  \item 一个铀-235 \omterm{def:g11:fundamental-interactions:ladder}{原子核}的质量（$A = 235$），再求每秒消耗的铀-235
        质量？
  \item 推算一年内消耗的铀-235 质量，以\omterm{def:g10:orders-of-magnitude:unit}{千克}计。
  \item 燃料富集至 $4.0\%$ 的铀-235：一年需要多少质量的燃料，
        一辆卡车能否装下？
\end{enumerate}

\medskip
\noindent\textbf{第二部分——煤的账本。}
\begin{enumerate}[resume]
  \item 计算该城市一年内消耗的\omterm{def:g10:energy-conservation:forms}{热能}。
  \item 需要多少质量的煤才能提供这些\omterm{def:g10:energy-conservation:forms}{热能}？
  \item 计算一年内（煤的质量）/（铀-235 的质量）之比，并
        对照\cref{ex:g12:nuclear-energy:ladder}中的阶梯加以检验。
  \item 一列货运火车运载 \qty{3.0e6}{kg}：一年需要多少列，
        每天需要多少列？用一句话描述这两条供应线。
\end{enumerate}

\medskip
\noindent\textbf{第三部分——\omterm{def:g12:nuclear-energy:fusion}{聚变}之梦。}
\begin{enumerate}[resume]
  \item 用\omterm{def:g11:work-of-force:work}{焦耳}表示一次 D--T \omterm{def:g12:nuclear-energy:fusion}{聚变}的\omterm{def:g10:energy-conservation:energy}{能量}。
  \item 计算每\omterm{def:g10:orders-of-magnitude:unit}{千克} D--T 燃料释放的\omterm{def:g10:energy-conservation:energy}{能量}，并与
        阶梯比较。
  \item 该城市一年所需的 D--T 燃料有多少质量？其中
        氘（\qty{5.029}{u} 中的 $2.014$）占多少质量？
  \item 含有该氘的海水体积是多少？与
        一座奥运标准游泳池（约 \qty{2.5e3}{m^3}）相比如何？（氚由锂增殖而来。）
  \item 这座电厂为何无须担心\omterm{def:g12:nuclear-energy:chain}{临界质量}与\omterm{def:g12:nuclear-energy:chain}{链式反应}？
        那么，真正的难题是什么？
\end{enumerate}

\medskip
\noindent\textbf{第四部分——以质量换取\omterm{def:g10:energy-conservation:energy}{能量}的恒星。}
\begin{enumerate}[resume]
  \item 由 $E = mc^2$，太阳每秒转化多少质量？
  \item 每秒、每年，\emph{任何}热\omterm{def:g11:work-of-force:power}{功率}为
        \qty{3.0}{GW} 的能源必须转化多少质量？验证三种燃料
        都损失同样的质量。
  \item 太阳在其 $\num{4.5e9}$
        年中总共辐射掉多少质量？这占太阳质量的几分之几？
  \item \omterm{def:g12:nuclear-energy:fusion}{聚变}可利用太阳质量的约 $0.07\%$（即其核心中可燃烧的
        部分）：估算它作为恒星的总寿命，以及剩余的质量。
  \item 用三句话结束这份核算：该城市一年所需的煤（以吨计）、
        铀-235（以\omterm{def:g10:orders-of-magnitude:unit}{千克}计）与 D--T 燃料（以\omterm{def:g10:orders-of-magnitude:unit}{千克}计），以及
        在每种方案中化为该城市光明的那一\omterm{def:g10:orders-of-magnitude:unit}{千克}质量。
\end{enumerate}
\end{problem}
